% Options for packages loaded elsewhere
% Options for packages loaded elsewhere
\PassOptionsToPackage{unicode}{hyperref}
\PassOptionsToPackage{hyphens}{url}
%
\documentclass[
  brazilian,
  11pt,
  a4paper,
  open = right,
  twoside = true]{scrbook}
\usepackage{xcolor}
\usepackage[top=2.5cm,bottom=2.5cm,inner=3cm,textwidth=12cm,marginparwidth=4.5cm,marginparsep=5mm,outer=6cm,twoside=true]{geometry}
\usepackage{amsmath,amssymb}
\setcounter{secnumdepth}{5}
\usepackage{iftex}
\ifPDFTeX
  \usepackage[T1]{fontenc}
  \usepackage[utf8]{inputenc}
  \usepackage{textcomp} % provide euro and other symbols
\else % if luatex or xetex
  \usepackage{unicode-math} % this also loads fontspec
  \defaultfontfeatures{Scale=MatchLowercase}
  \defaultfontfeatures[\rmfamily]{Ligatures=TeX,Scale=1}
\fi
\usepackage{lmodern}
\ifPDFTeX\else
  % xetex/luatex font selection
  \setmainfont[]{TeX Gyre Termes}
  \setsansfont[]{Fira Sans}
  \setmonofont[]{Fira Mono}
  \setmathfont[]{Fira Math}
\fi
% Use upquote if available, for straight quotes in verbatim environments
\IfFileExists{upquote.sty}{\usepackage{upquote}}{}
\IfFileExists{microtype.sty}{% use microtype if available
  \usepackage[]{microtype}
  \UseMicrotypeSet[protrusion]{basicmath} % disable protrusion for tt fonts
}{}
\makeatletter
\@ifundefined{KOMAClassName}{% if non-KOMA class
  \IfFileExists{parskip.sty}{%
    \usepackage{parskip}
  }{% else
    \setlength{\parindent}{0pt}
    \setlength{\parskip}{6pt plus 2pt minus 1pt}}
}{% if KOMA class
  \KOMAoptions{parskip=half}}
\makeatother
% Make \paragraph and \subparagraph free-standing
\makeatletter
\ifx\paragraph\undefined\else
  \let\oldparagraph\paragraph
  \renewcommand{\paragraph}{
    \@ifstar
      \xxxParagraphStar
      \xxxParagraphNoStar
  }
  \newcommand{\xxxParagraphStar}[1]{\oldparagraph*{#1}\mbox{}}
  \newcommand{\xxxParagraphNoStar}[1]{\oldparagraph{#1}\mbox{}}
\fi
\ifx\subparagraph\undefined\else
  \let\oldsubparagraph\subparagraph
  \renewcommand{\subparagraph}{
    \@ifstar
      \xxxSubParagraphStar
      \xxxSubParagraphNoStar
  }
  \newcommand{\xxxSubParagraphStar}[1]{\oldsubparagraph*{#1}\mbox{}}
  \newcommand{\xxxSubParagraphNoStar}[1]{\oldsubparagraph{#1}\mbox{}}
\fi
\makeatother

\usepackage{color}
\usepackage{fancyvrb}
\newcommand{\VerbBar}{|}
\newcommand{\VERB}{\Verb[commandchars=\\\{\}]}
\DefineVerbatimEnvironment{Highlighting}{Verbatim}{commandchars=\\\{\}}
% Add ',fontsize=\small' for more characters per line
\usepackage{framed}
\definecolor{shadecolor}{RGB}{253,251,247}
\newenvironment{Shaded}{\begin{snugshade}}{\end{snugshade}}
\newcommand{\AlertTok}[1]{\textcolor[rgb]{0.10,0.10,0.10}{#1}}
\newcommand{\AnnotationTok}[1]{\textcolor[rgb]{0.10,0.10,0.10}{#1}}
\newcommand{\AttributeTok}[1]{\textcolor[rgb]{0.00,0.28,0.67}{#1}}
\newcommand{\BaseNTok}[1]{\textcolor[rgb]{0.10,0.10,0.10}{#1}}
\newcommand{\BuiltInTok}[1]{\textcolor[rgb]{0.00,0.28,0.67}{#1}}
\newcommand{\CharTok}[1]{\textcolor[rgb]{0.95,0.43,0.13}{#1}}
\newcommand{\CommentTok}[1]{\textcolor[rgb]{0.32,0.36,0.40}{\textit{#1}}}
\newcommand{\CommentVarTok}[1]{\textcolor[rgb]{0.10,0.10,0.10}{#1}}
\newcommand{\ConstantTok}[1]{\textcolor[rgb]{0.00,0.77,0.85}{#1}}
\newcommand{\ControlFlowTok}[1]{\textcolor[rgb]{0.64,0.23,0.08}{\textbf{#1}}}
\newcommand{\DataTypeTok}[1]{\textcolor[rgb]{0.64,0.23,0.08}{#1}}
\newcommand{\DecValTok}[1]{\textcolor[rgb]{0.55,0.29,0.58}{#1}}
\newcommand{\DocumentationTok}[1]{\textcolor[rgb]{0.16,0.46,0.83}{#1}}
\newcommand{\ErrorTok}[1]{\textcolor[rgb]{0.95,0.43,0.13}{\underline{#1}}}
\newcommand{\ExtensionTok}[1]{\textcolor[rgb]{0.00,0.28,0.67}{\textbf{#1}}}
\newcommand{\FloatTok}[1]{\textcolor[rgb]{0.11,0.47,0.80}{#1}}
\newcommand{\FunctionTok}[1]{\textcolor[rgb]{1.00,0.23,0.00}{#1}}
\newcommand{\ImportTok}[1]{\textcolor[rgb]{0.64,0.23,0.08}{#1}}
\newcommand{\InformationTok}[1]{\textcolor[rgb]{0.16,0.46,0.83}{#1}}
\newcommand{\KeywordTok}[1]{\textcolor[rgb]{1.00,0.23,0.00}{\textbf{#1}}}
\newcommand{\NormalTok}[1]{\textcolor[rgb]{0.10,0.10,0.10}{#1}}
\newcommand{\OperatorTok}[1]{\textcolor[rgb]{0.10,0.10,0.10}{#1}}
\newcommand{\OtherTok}[1]{\textcolor[rgb]{0.10,0.10,0.10}{#1}}
\newcommand{\PreprocessorTok}[1]{\textcolor[rgb]{0.95,0.43,0.13}{#1}}
\newcommand{\RegionMarkerTok}[1]{\textcolor[rgb]{0.16,0.46,0.83}{#1}}
\newcommand{\SpecialCharTok}[1]{\textcolor[rgb]{0.85,0.00,0.31}{#1}}
\newcommand{\SpecialStringTok}[1]{\textcolor[rgb]{0.85,0.00,0.31}{#1}}
\newcommand{\StringTok}[1]{\textcolor[rgb]{0.00,0.72,0.25}{#1}}
\newcommand{\VariableTok}[1]{\textcolor[rgb]{0.10,0.10,0.10}{#1}}
\newcommand{\VerbatimStringTok}[1]{\textcolor[rgb]{0.10,0.10,0.10}{#1}}
\newcommand{\WarningTok}[1]{\textcolor[rgb]{0.40,0.22,0.33}{\textit{#1}}}

\usepackage{longtable,booktabs,array}
\newcounter{none} % for unnumbered tables
\usepackage{calc} % for calculating minipage widths
% Correct order of tables after \paragraph or \subparagraph
\usepackage{etoolbox}
\makeatletter
\patchcmd\longtable{\par}{\if@noskipsec\mbox{}\fi\par}{}{}
\makeatother
% Allow footnotes in longtable head/foot
\IfFileExists{footnotehyper.sty}{\usepackage{footnotehyper}}{\usepackage{footnote}}
\makesavenoteenv{longtable}
\usepackage{graphicx}
\makeatletter
\newsavebox\pandoc@box
\newcommand*\pandocbounded[1]{% scales image to fit in text height/width
  \sbox\pandoc@box{#1}%
  \Gscale@div\@tempa{\textheight}{\dimexpr\ht\pandoc@box+\dp\pandoc@box\relax}%
  \Gscale@div\@tempb{\linewidth}{\wd\pandoc@box}%
  \ifdim\@tempb\p@<\@tempa\p@\let\@tempa\@tempb\fi% select the smaller of both
  \ifdim\@tempa\p@<\p@\scalebox{\@tempa}{\usebox\pandoc@box}%
  \else\usebox{\pandoc@box}%
  \fi%
}
% Set default figure placement to htbp
\def\fps@figure{htbp}
\makeatother



\ifLuaTeX
\usepackage[bidi=basic,shorthands=off]{babel}
\else
\usepackage[bidi=default,shorthands=off]{babel}
\fi
\ifPDFTeX
\else
\babelfont{rm}[]{TeX Gyre Termes}
\fi
\ifLuaTeX
  \usepackage{selnolig} % disable illegal ligatures
\fi


\setlength{\emergencystretch}{3em} % prevent overfull lines

\providecommand{\tightlist}{%
  \setlength{\itemsep}{0pt}\setlength{\parskip}{0pt}}



 


% =========================
% TIPOGRAFIA
% =========================

\usepackage{microtype}
\usepackage{lettrine}

% =========================
% CORES INFOSOC (Cubismo laranja creme azul)
% =========================

\usepackage{xcolor}
\usepackage{pagecolor}
\usepackage{afterpage}
\usepackage{etoolbox}

% Definição exata da paleta de apoio
\definecolor{laranjaqueimado}{HTML}{A33B14}
\definecolor{laranjavibrante}{HTML}{FF3B00}
\definecolor{laranja}{HTML}{F26D21}
\definecolor{pessego}{HTML}{FFC299}
\definecolor{creme}{HTML}{FDFBF7}
\definecolor{bege}{HTML}{D9CBB8}
\definecolor{azul}{HTML}{2A75D3}
\definecolor{ciano}{HTML}{00C4D9}
\definecolor{azulcobalto}{HTML}{0047AB}
\definecolor{textoescuro}{HTML}{1A1A1A}

% =========================
% KOMA - FONTES E CORES ESTRUTURAIS
% =========================

%Todos os títulos, incluindo o sumário
\setkomafont{disposition}{\sffamily\color{azulcobalto}}

% Título da primeira página do livro
\setkomafont{title}{\sffamily\color{laranjavibrante}}
\setkomafont{author}{\sffamily}
\setkomafont{date}{\sffamily}

%Numeração da parte
\setkomafont{partnumber}{\sffamily\bfseries\color{laranjaqueimado}\Huge}
\setkomafont{part}{\sffamily\bfseries\color{laranjavibrante}\fontsize{60}{76}\selectfont}

% Capítulo é definido mais abaixos

% Seções
\setkomafont{section}{\sffamily\bfseries\color{laranjavibrante}\Large}
\setkomafont{subsection}{\sffamily\color{azul}}
\setkomafont{subsubsection}{\sffmaily\color{azulcobalto}}

% =========================
% TOC
% =========================

\usepackage{tocbasic}

% Comando auxiliar que recebe o texto da entrada como #1

\newcommand{\formatacaoparte}[1]{%
  \setlength{\fboxsep}{6pt}%
  \colorbox{pessego}{%
    \makebox[\dimexpr\linewidth-2\fboxsep+28pt\relax][l]{%
      \color{laranjaqueimado}\sffamily\bfseries\large #1%
    }%
  }%
}

\DeclareTOCStyleEntry[
  entryformat=\formatacaoparte,
  pagenumberformat={\color{laranjaqueimado}\sffamily\bfseries},
  linefill={\leavevmode\color{laranjaqueimado}\TOCLineLeaderFill},
  beforeskip=1.5em
]{tocline}{part}


\DeclareTOCStyleEntry[
  entryformat={\color{laranjaqueimado}\sffamily},
  pagenumberformat={\color{laranjaqueimado}\sffamily\bfseries},
  linefill=\TOCLineLeaderFill
]{tocline}{chapter}

\DeclareTOCStyleEntry[
  entryformat={\color{laranjavibrante}\sffamily},
  pagenumberformat=\sffamily,
  linefill=\TOCLineLeaderFill
]{tocline}{section}

\DeclareTOCStyleEntry[
  entryformat={\color{azul}\sffamily},
  pagenumberformat=\sffamily
]{tocline}{subsection}

% =========================
% CAPÍTULOS
% =========================

\RedeclareSectionCommand[
  beforeskip=2.0cm,
  afterskip=1.0cm
]{chapter}

% Remove o ponto após o número (padrão KOMA é "1.")
\renewcommand*{\chapterformat}{\thechapter}

% #2 = número formatado | #3 = título
% \ifblank: se #2 vazio (cap. não numerado), não exibe a caixa
\renewcommand*{\chapterlinesformat}[3]{%
  \noindent
  \ifblank{#2}{}{%
    {%
      \setlength{\fboxsep}{0.45em}%
      \colorbox{bege}{%
        \hspace{0.3em}%
        {\color{white}\bfseries\huge%
          \rule[-0.25em]{0pt}{1.2em}#2}%
        \hspace{0.3em}%
      }%
    }%
    \par\vspace{0.7em}%
  }%
  {\sffamily\bfseries\Huge\color{laranjaqueimado}#3}%
  \par
}

% =========================
% PARTES
% =========================

\renewcommand*{\partformat}{\partname~\thepart}
\renewcommand*{\partpagestyle}{empty}

% \pretocmd: ativa a cor de fundo ANTES da página de parte ser composta.
% Utiliza o pêssego da nova paleta para um fundo de transição suave.
\pretocmd{\part}{%
  \cleardoublepage
  \pagecolor{pessego}%
}{}{}

% \partheadendvskip: chamado DENTRO da página de parte, ao final.
% \afterpage{\nopagecolor} é acionado após a página ser finalizada,
% restaurando o fundo branco puro para as páginas seguintes.
\renewcommand*{\partheadstartvskip}{\null\vfill}
\renewcommand*{\partheadendvskip}{%
  \vfill\null
  \afterpage{\nopagecolor}%
}

% =========================
% CABEÇALHO
% =========================

\usepackage{scrlayer-scrpage}
\clearpairofpagestyles

\automark[section]{chapter}

\ihead{\headmark}
\ohead{\pagemark}

% Elementos de navegação em azul cobalto garantem sobriedade e leitura eficiente
\setkomafont{pagehead}{\sffamily\small\color{laranja}}
\setkomafont{pagenumber}{\sffamily\small\color{laranja}}

\KOMAoptions{headsepline=true}
% Colorir a linha do cabeçalho
\addtokomafont{headsepline}{\color{bege}}

% =========================
% LEGENDAS (Figuras e Tabelas)
% =========================

% Altera a cor e a fonte do texto da legenda
\setkomafont{caption}{\sffamily\color{laranjaqueimado}}
\setcapindent{0em}
% Altera a cor e a fonte do rótulo ("Figura 1:", "Tabela 1:")
\setkomafont{captionlabel}{\sffamily\color{laranjaqueimado}}


% =========================
% NOTAS DE RODAPÉ
% =========================

% Altera a cor do texto da nota de rodapé
\setkomafont{footnote}{\sffamily\\color{azulcobalto}}

% Altera a cor do número/marcação da nota no rodapé da página
\setkomafont{footnotelabel}{\sffamily\\color{azulcobalto}}

% Opcional: Altera a cor do número sobrescrito no meio do texto principal
% \setkomafont{footnotereference}{\color{azulcobalto}}

% =========================
% FONTE DO CORPO DAS TABELAS
% =========================

% Injeta \sffamily no início dos principais ambientes de tabela
\AtBeginEnvironment{table}{\sffamily}
\AtBeginEnvironment{tabular}{\sffamily}
\AtBeginEnvironment{tabular*}{\sffamily}
\AtBeginEnvironment{longtable}{\sffamily}

% =========================
% CONTROLE DE FONTE EM MINIPAGE (TABELAS)
% =========================

% 1. Cria a variável de estado lógico (inicia como falsa por padrão)
\newtoggle{dentrodetabela}

% 2. Altera o estado para verdadeiro ao iniciar um ambiente de tabela
\AtBeginEnvironment{table}{\toggletrue{dentrodetabela}}

% 3. Injeta a verificação condicional no início de toda minipage
\AtBeginEnvironment{minipage}{%
  \iftoggle{dentrodetabela}{\sffamily\small}{}%
}

% =========================
% NOTAS MARGINAIS (Integração Quarto)
% =========================

% 1. Assume o controle sobre as notas geradas por marcações do Quarto (.aside)
\usepackage{marginnote}
\renewcommand*{\marginfont}{\small\color{azulcobalto}}

% 2. Mantém a intercepção do comando KOMA-Script para marcações manuais no texto
\let\oldmarginpar\marginpar
\renewcommand{\marginpar}[1]{\oldmarginpar{\small\color{azulcobalto}#1}}

\let\oldmarginline\marginline
\renewcommand{\marginline}[1]{\oldmarginline{\small\color{azulcobalto}#1}}
\usepackage{booktabs}
\usepackage{caption}
\usepackage{longtable}
\usepackage{colortbl}
\usepackage{array}
\usepackage{anyfontsize}
\usepackage{multirow}
\makeatletter
\@ifpackageloaded{tcolorbox}{}{\usepackage[skins,breakable]{tcolorbox}}
\@ifpackageloaded{fontawesome5}{}{\usepackage{fontawesome5}}
\definecolor{quarto-callout-color}{HTML}{909090}
\definecolor{quarto-callout-note-color}{HTML}{0758E5}
\definecolor{quarto-callout-important-color}{HTML}{CC1914}
\definecolor{quarto-callout-warning-color}{HTML}{EB9113}
\definecolor{quarto-callout-tip-color}{HTML}{00A047}
\definecolor{quarto-callout-caution-color}{HTML}{FC5300}
\definecolor{quarto-callout-color-frame}{HTML}{acacac}
\definecolor{quarto-callout-note-color-frame}{HTML}{4582ec}
\definecolor{quarto-callout-important-color-frame}{HTML}{d9534f}
\definecolor{quarto-callout-warning-color-frame}{HTML}{f0ad4e}
\definecolor{quarto-callout-tip-color-frame}{HTML}{02b875}
\definecolor{quarto-callout-caution-color-frame}{HTML}{fd7e14}
\makeatother
\makeatletter
\@ifpackageloaded{bookmark}{}{\usepackage{bookmark}}
\makeatother
\makeatletter
\@ifpackageloaded{caption}{}{\usepackage{caption}}
\AtBeginDocument{%
\ifdefined\contentsname
  \renewcommand*\contentsname{Índice}
\else
  \newcommand\contentsname{Índice}
\fi
\ifdefined\listfigurename
  \renewcommand*\listfigurename{Lista de Figuras}
\else
  \newcommand\listfigurename{Lista de Figuras}
\fi
\ifdefined\listtablename
  \renewcommand*\listtablename{Lista de Tabelas}
\else
  \newcommand\listtablename{Lista de Tabelas}
\fi
\ifdefined\figurename
  \renewcommand*\figurename{Figura}
\else
  \newcommand\figurename{Figura}
\fi
\ifdefined\tablename
  \renewcommand*\tablename{Tabela}
\else
  \newcommand\tablename{Tabela}
\fi
}
\@ifpackageloaded{float}{}{\usepackage{float}}
\floatstyle{ruled}
\@ifundefined{c@chapter}{\newfloat{codelisting}{h}{lop}}{\newfloat{codelisting}{h}{lop}[chapter]}
\floatname{codelisting}{Listagem}
\newcommand*\listoflistings{\listof{codelisting}{Lista de Listagens}}
\makeatother
\makeatletter
\makeatother
\makeatletter
\@ifpackageloaded{caption}{}{\usepackage{caption}}
\@ifpackageloaded{subcaption}{}{\usepackage{subcaption}}
\makeatother
\makeatletter
\@ifpackageloaded{sidenotes}{}{\usepackage{sidenotes}}
\@ifpackageloaded{marginnote}{}{\usepackage{marginnote}}
\makeatother
\newcounter{quartocallouttipno}
\newcommand{\quartocallouttip}[1]{\refstepcounter{quartocallouttipno}\label{#1}}
\usepackage{bookmark}
\IfFileExists{xurl.sty}{\usepackage{xurl}}{} % add URL line breaks if available
\urlstyle{same}
\hypersetup{
  pdftitle={Métodos Computacionais para Indicadores Sociais e Políticas Públicas},
  pdfauthor={Ronaldo Baltar e Cláudia Baltar},
  pdflang={pt-BR},
  hidelinks,
  pdfcreator={LaTeX via pandoc}}


\title{Métodos Computacionais para Indicadores Sociais e Políticas
Públicas}
\usepackage{etoolbox}
\makeatletter
\providecommand{\subtitle}[1]{% add subtitle to \maketitle
  \apptocmd{\@title}{\par {\large #1 \par}}{}{}
}
\makeatother
\subtitle{Usando R}
\author{Ronaldo Baltar e Cláudia Baltar}
\date{2026}
\begin{document}
\frontmatter
\maketitle

\renewcommand*\contentsname{Índice}
{
\setcounter{tocdepth}{2}
\tableofcontents
}

\mainmatter
\bookmarksetup{startatroot}

\chapter*{Prefácio A}\label{prefuxe1cio-a}
\addcontentsline{toc}{chapter}{Prefácio A}

\markboth{Prefácio A}{Prefácio A}

Este é o texto do seu prefácio. Note que o comando no seu YAML garantiu
que a contagem começasse na Falsa Folha de Rosto, mas o impediu que o
número fosse impresso nas páginas iniciais.

\part{Fundamentos}

\chapter*{Introdução}\label{introduuxe7uxe3o}
\addcontentsline{toc}{chapter}{Introdução}

\markboth{Introdução}{Introdução}

Este conteúdo digital tem origem no curso online \emph{``Introdução aos
Indicadores Sociais com o software R''}, de autoria do Prof.~Ronaldo
Baltar e da Prof.ª Cláudia Siqueira Baltar. Ele foi desenvolvido como
uma atividade do Projeto de extensão ``Indicadores sociais como subsídio
para o monitoramento e avaliação das ações dos municípios paranaenses em
direção à Agenda 2030'', vinculado ao Observatório de Populações e
Políticas Públicas (ObPPP) e ao projeto de Informática aplicada à
Pesquisa Social (InfoSoc) da Universidade Estadual de Londrina (UEL).

\section*{Objetivo}\label{objetivo}
\addcontentsline{toc}{section}{Objetivo}

\markright{Objetivo}

O nosso foco inicial é a instalação e os primeiros passos com o
\textbf{R}, um software gratuito voltado para a análise de dados, e com
o \textbf{RStudio}, um ambiente integrado de desenvolvimento (IDE)
também gratuito. O RStudio oferece diversos recursos práticos para
facilitar a escrita, a elaboração e a visualização dos códigos que você
criará utilizando o R.

\section*{Preparando o Ambiente de
Trabalho}\label{preparando-o-ambiente-de-trabalho}
\addcontentsline{toc}{section}{Preparando o Ambiente de Trabalho}

\markright{Preparando o Ambiente de Trabalho}

\subsection*{Instalando o R}\label{instalando-o-r}
\addcontentsline{toc}{subsection}{Instalando o R}

Para realizar a instalação da linguagem R, você deve acessar o site do
repositório padrão \textbf{CRAN} através do link abaixo:

\url{https://cran.r-project.org}

Ao acessar o site, localize o \emph{link} da versão compatível com o seu
sistema operacional (Windows, Linux ou MacOS). Baixe o arquivo de
instalação mais recente e execute-o diretamente no seu computador.

\subsection*{Instalando o RStudio
(Posit)}\label{instalando-o-rstudio-posit}
\addcontentsline{toc}{subsection}{Instalando o RStudio (Posit)}

O RStudio é distribuído em duas versões: \emph{Desktop} (desenvolvida
para uso local em computadores individuais) e \emph{Server} (servidor).
Ambas estão disponíveis sob licença comercial ou gratuita (\emph{open
source}) para Linux, Mac e Windows.

Em novembro de 2022, o RStudio mudou sua marca para \textbf{Posit}, mas
os procedimentos de uso e instalação continuam os mesmos. Lembre-se de
duas regras fundamentais antes da instalação: 1. O RStudio aceita apenas
instalação em \textbf{computadores de 64-bit}. 2. Você deve se
certificar de que \textbf{o R já tenha sido instalado corretamente}
antes de prosseguir com o RStudio.

O arquivo de instalação da versão Desktop pode ser baixado no seguinte
endereço: \url{https://posit.co/download/rstudio-desktop/}

Se tiver dificuldades ao longo do processo, você pode acompanhar este
\href{https://www.iorad.com/player/1815363/Como-instalar-o-R-e-o-RStudio}{Tutorial
Interativo de Instalação do R e RStudio}.

\begin{tcolorbox}[enhanced jigsaw, arc=.35mm, bottomrule=.15mm, bottomtitle=1mm, breakable, colback=white, colbacktitle=quarto-callout-note-color!10!white, colframe=quarto-callout-note-color-frame, coltitle=black, left=2mm, leftrule=.75mm, opacityback=0, opacitybacktitle=0.6, rightrule=.15mm, title=\textcolor{quarto-callout-note-color}{\faInfo}\hspace{0.5em}{Sobre a Instalação do RTools (Apenas para Windows)}, titlerule=0mm, toprule=.15mm, toptitle=1mm]

Ao instalar alguns pacotes que utilizam recursos computacionais mais
avançados (como o \texttt{ggplot2}), eles podem solicitar compilação, ou
seja, conversão de linguagens como C/C++ ou Fortran. Nesses casos, o R
pode alertar que existe uma versão em código fonte mais recente do que a
versão binária pronta para instalação.

\textbf{Não há nenhum problema em ignorar essa mensagem} e continuar
utilizando a versão binária disponível.

No entanto, se você deseja compilar códigos fontes no seu próprio
computador e utiliza o sistema Windows, precisará instalar uma
ferramenta de suporte chamada \textbf{RTools}. O RTools ocupa cerca de
\textbf{1 GB de espaço no disco rígido} e pode ser baixado em:
https://cran.r-project.org/bin/windows/Rtools/

\end{tcolorbox}

\chapter{Texto modelo}\label{texto-modelo}

\lettrine{A}{ntes} iria o tableofcontents, mas não vai mais. Vamos ver.
Lorem class magna sapien duis convallis, convallis bibendum aenean quam.
Enim porta risus arcu turpis lobortis nisl risus maecenas cubilia
lobortis! Neque egestas condimentum. Adipiscing libero nisi neque proin
libero ac ornare; mauris egestas cursus! Condimentum conubia phasellus
class molestie vehicula, nisl semper nec massa.

Litora class quisque ut porttitor nostra at fames sollicitudin. Vel
mauris penatibus eu magna facilisis accumsan eleifend lacinia! Leo mi
vehicula curabitur accumsan ultricies pulvinar placerat? Vitae erat,
viverra hendrerit curabitur id -- praesent aliquet aenean vestibulum
praesent mauris. In luctus ridiculus quisque neque auctor vel quis
dictumst pretium molestie laoreet. A nisl?

\begin{Shaded}
\begin{Highlighting}[numbers=left,,]
\FunctionTok{library}\NormalTok{(ggplot2)}
\FunctionTok{source}\NormalTok{(}\StringTok{"tema\_livro.R"}\NormalTok{)}
\FunctionTok{definir\_tema\_livro}\NormalTok{() }

\NormalTok{mtcars2 }\OtherTok{\textless{}{-}}\NormalTok{ mtcars}
\NormalTok{mtcars2}\SpecialCharTok{$}\NormalTok{am }\OtherTok{\textless{}{-}} \FunctionTok{factor}\NormalTok{(}
\NormalTok{  mtcars}\SpecialCharTok{$}\NormalTok{am, }\AttributeTok{labels =} \FunctionTok{c}\NormalTok{(}\StringTok{\textquotesingle{}automatic\textquotesingle{}}\NormalTok{, }\StringTok{\textquotesingle{}manual\textquotesingle{}}\NormalTok{)}
\NormalTok{)}
\FunctionTok{ggplot}\NormalTok{(mtcars2, }\FunctionTok{aes}\NormalTok{(hp, mpg, }\AttributeTok{color =}\NormalTok{ am)) }\SpecialCharTok{+}
  \FunctionTok{geom\_point}\NormalTok{() }\SpecialCharTok{+}
  \FunctionTok{geom\_smooth}\NormalTok{(}\AttributeTok{formula =}\NormalTok{ y }\SpecialCharTok{\textasciitilde{}}\NormalTok{ x, }\AttributeTok{method =} \StringTok{"loess"}\NormalTok{) }\SpecialCharTok{+}
  \FunctionTok{theme}\NormalTok{(}\AttributeTok{legend.position =} \StringTok{\textquotesingle{}bottom\textquotesingle{}}\NormalTok{) }
\end{Highlighting}
\end{Shaded}

\begin{marginfigure}

\centering{

\pandocbounded{\includegraphics[keepaspectratio]{modelo_files/figure-pdf/fig-mtcars-1.png}}

}

\caption{\label{fig-mtcars}MPG vs horsepower, colored by transmission.}

\end{marginfigure}%

Lorem semper aliquam felis enim gravida velit himenaeos purus dis! Ad
fames massa sodales nisi odio non conubia. Mi ligula pellentesque
ultrices tristique, turpis mi, metus fringilla. Fames iaculis natoque
ullamcorper aliquet, vitae inceptos suscipit risus pharetra nascetur
leo! Faucibus ad duis semper vehicula condimentum in massa? Nisl
facilisis id magnis convallis varius imperdiet libero duis. Iaculis
dictumst cum vivamus mollis mattis felis sociis. Consequat viverra
tristique eleifend accumsan volutpat, arcu, sapien vel fringilla.

Amet sapien ultrices ridiculus ridiculus libero luctus. Posuere purus
varius proin, ullamcorper auctor natoque imperdiet urna pretium
vestibulum accumsan, habitasse quam. Pharetra venenatis ridiculus ligula
quis, purus sed litora platea viverra imperdiet diam. Blandit elementum
odio id, feugiat cubilia lacinia. Faucibus mus conubia condimentum
ultricies leo fusce netus sem venenatis nullam.

You an also place tables in the margin of your document by specifying
column: margin.

\begin{Shaded}
\begin{Highlighting}[numbers=left,,]
\NormalTok{knitr}\SpecialCharTok{::}\FunctionTok{kable}\NormalTok{(}
\NormalTok{  mtcars[}\DecValTok{1}\SpecialCharTok{:}\DecValTok{3}\NormalTok{, }\DecValTok{1}\SpecialCharTok{:}\DecValTok{3}\NormalTok{]}
\NormalTok{)}
\end{Highlighting}
\end{Shaded}

{\def\LTcaptype{none} % do not increment counter
\begin{longtable}[]{@{}lrrr@{}}
\toprule\noalign{}
& mpg & cyl & disp \\
\midrule\noalign{}
\endhead
\bottomrule\noalign{}
\endlastfoot
Mazda RX4 & 21.0 & 6 & 160 \\
Mazda RX4 Wag & 21.0 & 6 & 160 \\
Datsun 710 & 22.8 & 4 & 108 \\
\end{longtable}
}

Elit magnis aptent morbi aptent conubia cras, ullamcorper velit! Per
curabitur vivamus fringilla, turpis mus quam cubilia blandit cras nisi.
Placerat facilisis vitae magnis orci urna porta sollicitudin penatibus.
Feugiat cursus congue laoreet, malesuada felis, imperdiet justo
hendrerit. Massa congue felis nam semper blandit ad porttitor interdum
leo.

You can also place content in the margin by targeting the margin column
using a div with the .column-margin class. For example:

\marginnote{\begin{footnotesize}

We know from \emph{the first fundamental theorem of calculus} that for
\(x\) in \([a, b]\):

\[\frac{d}{dx}\left( \int_{a}^{x} f(u)\,du\right)=f(x).\]

\end{footnotesize}}

Consectetur mauris quisque primis, facilisis imperdiet curae posuere.
Pellentesque a varius velit nibh, rutrum semper cum porta pellentesque!
Diam conubia magna; nisi vivamus, ultrices magnis a ac eleifend hac
lobortis congue! Fusce vivamus condimentum sociis iaculis magna luctus
neque. In ac est; eget tristique orci cursus rutrum scelerisque.

\begin{Shaded}
\begin{Highlighting}[numbers=left,,]
\NormalTok{mtcars[}\DecValTok{1}\SpecialCharTok{:}\DecValTok{6}\NormalTok{, }\DecValTok{1}\SpecialCharTok{:}\DecValTok{5}\NormalTok{] }\SpecialCharTok{|\textgreater{}}
  \FunctionTok{gt}\NormalTok{(}\AttributeTok{rownames\_to\_stub =} \ConstantTok{TRUE}\NormalTok{) }\SpecialCharTok{|\textgreater{}}
  \FunctionTok{tab\_header}\NormalTok{(}
    \AttributeTok{title =} \StringTok{"Tabela 1 {-} Desempenho de Veículos Selecionados"}\NormalTok{,}
    \AttributeTok{subtitle =} \StringTok{"Consumo, número de cilindros e deslocamento volumétrico"}
\NormalTok{  ) }\SpecialCharTok{|\textgreater{}}
  \FunctionTok{tab\_source\_note}\NormalTok{(}
    \AttributeTok{source\_note =} \StringTok{"Fonte: Motor Trend Car Road Tests (1974)."}
\NormalTok{  ) }\SpecialCharTok{|\textgreater{}}
  \FunctionTok{tab\_footnote}\NormalTok{(}
    \AttributeTok{footnote =} \StringTok{"Milhas por galão norte{-}americano."}\NormalTok{,}
    \AttributeTok{locations =} \FunctionTok{cells\_column\_labels}\NormalTok{(}\AttributeTok{columns =}\NormalTok{ mpg)}
\NormalTok{  ) }\SpecialCharTok{|\textgreater{}}
  \FunctionTok{tema\_tabela\_livro}\NormalTok{()}
\end{Highlighting}
\end{Shaded}

\begin{table}
\caption*{
{\fontsize{10}{13}\selectfont  Tabela 1 - Desempenho de Veículos Selecionados\fontsize{9}{11}\selectfont } \\ 
{\fontsize{9}{11}\selectfont  Consumo, número de cilindros e deslocamento volumétrico\fontsize{9}{11}\selectfont }
} 
\fontsize{9.0pt}{11.0pt}\selectfont
\begin{tabular*}{\linewidth}{@{\extracolsep{\fill}}l|rrrrr}
\toprule
 & mpg\textsuperscript{\textit{1}} & cyl & disp & hp & drat \\ 
\midrule\addlinespace[2.5pt]
Mazda RX4 & 21.0 & 6 & 160 & 110 & 3.90 \\ 
Mazda RX4 Wag & 21.0 & 6 & 160 & 110 & 3.90 \\ 
Datsun 710 & 22.8 & 4 & 108 & 93 & 3.85 \\ 
Hornet 4 Drive & 21.4 & 6 & 258 & 110 & 3.08 \\ 
Hornet Sportabout & 18.7 & 8 & 360 & 175 & 3.15 \\ 
Valiant & 18.1 & 6 & 225 & 105 & 2.76 \\ 
\bottomrule
\end{tabular*}
\begin{minipage}{\linewidth}
\vspace{.05em}
{\fontsize{9}{11}\selectfont \textcolor[HTML]{0047AB}{\textsuperscript{\textit{1}}Milhas por galão norte-americano.}}\\
{\fontsize{9}{11}\selectfont \textcolor[HTML]{0047AB}{Fonte: Motor Trend Car Road Tests (1974).\\
}}\end{minipage}
\end{table}

Consectetur ac quam penatibus ac arcu, imperdiet consequat scelerisque
nisl commodo nisi? Magna dictumst quam purus, cursus egestas luctus
posuere hendrerit. Facilisis auctor platea, tempus interdum nunc natoque
senectus tempus, justo diam etiam malesuada dui? Cursus eros eget
facilisis class, porta sapien praesent dictumst. Elementum rutrum rutrum
ut auctor ante felis: maecenas posuere vivamus lectus leo? Natoque odio!

Asides allow you to place content aside from the content it is placed
in. Asides look like footnotes, but do not include the footnote mark
(the superscript number).

{\marginnote{\begin{footnotesize}This is a span that has the class aside
which places it in the margin without a footnote
number.\end{footnotesize}}}

Consectetur sociis leo neque placerat turpis conubia non rutrum magnis
ligula vulputate nisl dignissim. Diam id libero laoreet, gravida fames
pellentesque.

For figures and tables, you may leave the content in the body of the
document while placing the caption in the margin of the document. Using
cap-location: margin in a code cell or document front matter to control
this. For example:

\begin{Shaded}
\begin{Highlighting}[numbers=left,,]
\FunctionTok{library}\NormalTok{(ggplot2)}
\NormalTok{mtcars2 }\OtherTok{\textless{}{-}}\NormalTok{ mtcars}
\NormalTok{mtcars2}\SpecialCharTok{$}\NormalTok{am }\OtherTok{\textless{}{-}} \FunctionTok{factor}\NormalTok{(}
\NormalTok{  mtcars}\SpecialCharTok{$}\NormalTok{am, }\AttributeTok{labels =} \FunctionTok{c}\NormalTok{(}\StringTok{\textquotesingle{}automatic\textquotesingle{}}\NormalTok{, }\StringTok{\textquotesingle{}manual\textquotesingle{}}\NormalTok{)}
\NormalTok{)}
\FunctionTok{ggplot}\NormalTok{(mtcars2, }\FunctionTok{aes}\NormalTok{(hp, mpg, }\AttributeTok{color =}\NormalTok{ am)) }\SpecialCharTok{+}
  \FunctionTok{geom\_smooth}\NormalTok{(}\AttributeTok{formula =}\NormalTok{ y }\SpecialCharTok{\textasciitilde{}}\NormalTok{ x, }\AttributeTok{method =} \StringTok{"loess"}\NormalTok{) }\SpecialCharTok{+}
  \FunctionTok{geom\_point}\NormalTok{() }\SpecialCharTok{+}
  \FunctionTok{theme}\NormalTok{(}\AttributeTok{legend.position =} \StringTok{\textquotesingle{}bottom\textquotesingle{}}\NormalTok{) }
\end{Highlighting}
\end{Shaded}

\begin{figure}[H]

{\centering \pandocbounded{\includegraphics[keepaspectratio]{modelo_files/figure-pdf/unnamed-chunk-3-1.png}}

}

\caption{MPG vs horsepower, colored by transmission.}

\end{figure}%

Consectetur blandit pellentesque arcu praesent tortor cras enim.
Pellentesque sem suspendisse dis, vulputate egestas nulla at commodo
fermentum ante, curabitur rutrum? Feugiat ultrices fermentum tristique
interdum tincidunt eleifend tempus massa turpis mollis. Metus arcu
magnis mauris vitae fermentum, arcu magna orci sem. Nulla torquent
dictum proin.

Amet aliquet rhoncus volutpat auctor placerat lectus, curabitur:
venenatis non ultricies. Condimentum ac vitae netus, massa fringilla
volutpat nascetur, risus nisi. Hendrerit lobortis libero massa senectus
tristique, sodales nullam potenti mi quam? Aliquet potenti imperdiet
pharetra volutpat taciti arcu. Hac commodo non duis in curae: hendrerit
suscipit, netus blandit vivamus venenatis sapien fusce dapibus sagittis.

For computational output, you can specify the page column in your code
cell options. For example:

\begin{Shaded}
\begin{Highlighting}[numbers=left,,]
\NormalTok{knitr}\SpecialCharTok{::}\FunctionTok{kable}\NormalTok{(}
\NormalTok{  mtcars[}\DecValTok{1}\SpecialCharTok{:}\DecValTok{6}\NormalTok{, }\DecValTok{1}\SpecialCharTok{:}\DecValTok{10}\NormalTok{]}
\NormalTok{)}
\end{Highlighting}
\end{Shaded}

\begin{table*}

\begin{tabular}{lrrrrrrrrrr}
\toprule
 & mpg & cyl & disp & hp & drat & wt & qsec & vs & am & gear\\
\midrule
Mazda RX4 & 21.0 & 6 & 160 & 110 & 3.90 & 2.620 & 16.46 & 0 & 1 & 4\\
Mazda RX4
Wag & 21.0 & 6 & 160 & 110 & 3.90 & 2.875 & 17.02 & 0 & 1 & 4\\
Datsun 710 & 22.8 & 4 & 108 & 93 & 3.85 & 2.320 & 18.61 & 1 & 1 & 4\\
Hornet 4
Drive & 21.4 & 6 & 258 & 110 & 3.08 & 3.215 & 19.44 & 1 & 0 & 3\\
Hornet
Sportabout & 18.7 & 8 & 360 & 175 & 3.15 & 3.440 & 17.02 & 0 & 0 & 3\\
Valiant & 18.1 & 6 & 225 & 105 & 2.76 & 3.460 & 20.22 & 1 & 0 & 3\\
\bottomrule
\end{tabular}

\end{table*}%

Ipsum ligula a, euismod habitasse imperdiet dui dui purus dictumst. Sem
vitae sapien pharetra at lacus fringilla praesent, risus gravida
faucibus. At orci dictum consequat aliquet tempus aenean, cubilia,
ultricies potenti quisque sapien lectus quis.

\begin{tcolorbox}[enhanced jigsaw, arc=.35mm, bottomrule=.15mm, breakable, colback=white, colframe=quarto-callout-note-color-frame, left=2mm, leftrule=.75mm, opacityback=0, rightrule=.15mm, toprule=.15mm]
\begin{minipage}[t]{5.5mm}
\textcolor{quarto-callout-note-color}{\faInfo}
\end{minipage}%
\begin{minipage}[t]{\textwidth - 5.5mm}

\vspace{-3mm}\textbf{Pay Attention}\vspace{3mm}

Using callouts is an effective way to highlight content that your reader
give special consideration or attention.

\end{minipage}%
\end{tcolorbox}

Adipiscing mollis ante dapibus, est pretium auctor molestie hac, natoque
aliquet dis cum. Dictumst dictum rutrum, habitant condimentum dignissim
vehicula semper, tristique habitasse ornare volutpat! Ante dis varius
faucibus rhoncus cubilia odio nisl. Tincidunt orci dui ultrices ad
sodales, aenean et laoreet. Vehicula fringilla, ut lectus fringilla,
dignissim nam cubilia convallis elementum. Habitasse nisi integer, vitae
mi aliquam molestie eleifend at.

To cross-reference a callout, add an ID attribute that starts with the
appropriate callout prefix (see Table 1). You can then reference the
callout using the usual @ syntax. For example, here we add the ID
\#tip-example to the callout, and then refer back to it:

\begin{tcolorbox}[enhanced jigsaw, arc=.35mm, bottomrule=.15mm, bottomtitle=1mm, breakable, colback=white, colbacktitle=quarto-callout-tip-color!10!white, colframe=quarto-callout-tip-color-frame, coltitle=black, left=2mm, leftrule=.75mm, opacityback=0, opacitybacktitle=0.6, rightrule=.15mm, title=\textcolor{quarto-callout-tip-color}{\faLightbulb}\hspace{0.5em}{Dica \ref*{tip-example}: Cross-Referencing a Tip}, titlerule=0mm, toprule=.15mm, toptitle=1mm]

\quartocallouttip{tip-example} 

Add an ID starting with \texttt{\#tip-} to reference a tip.

\end{tcolorbox}

See Tip~\ref{tip-example}\ldots{}

This renders as follows:

Lorem porttitor phasellus sem eget, dictum pellentesque iaculis ligula
pellentesque litora! Netus lobortis id faucibus aenean etiam faucibus
gravida quis potenti accumsan. Mus litora molestie, pulvinar tempus
proin pharetra!

Consectetur placerat sollicitudin vivamus feugiat purus conubia -- eu
arcu nec taciti. Lectus sociis, tortor etiam, arcu dui; venenatis rutrum
vivamus! Torquent magna sapien; in magnis, morbi in, habitant quisque
interdum mauris. Ultricies fusce primis laoreet taciti sollicitudin
penatibus platea montes, primis interdum. Sagittis laoreet lacinia;
integer nullam purus etiam fusce!

\section{Sublorem}\label{sublorem}

Sit varius leo bibendum ante: elementum nisi per -- luctus aliquet
morbi. Bibendum nam nibh: ante: placerat mauris, urna habitasse ante
ornare penatibus eros. Suscipit taciti mus proin -- porta cum egestas
pharetra integer hendrerit? Tristique vehicula rutrum: vitae proin
semper magnis?

Elit nisi class\index{class} parturient sodales! Ultrices cubilia morbi
condimentum mus molestie habitasse nostra lacus molestie. Ad vitae
litora metus euismod suspendisse at sodales mi urna taciti pulvinar.
Bibendum ante dignissim netus donec commodo -- orci dis morbi tempor
parturient?

\subsection{Sub sub lorem}\label{sub-sub-lorem}

Consectetur curabitur et cubilia scelerisque id enim id aliquet, lacinia
venenatis justo facilisi? Inceptos netus taciti sodales turpis habitant:
nam himenaeos cras magna. Nullam parturient mus arcu platea sodales
pulvinar. Venenatis sem eros facilisis conubia aenean; aenean netus
facilisi. Scelerisque eros mauris libero hac consequat facilisi ornare
non sociis orci sollicitudin! Commodo hac aliquam duis, fermentum quis
cursus odio lacus montes odio commodo, tortor natoque consequat mi
eleifend porta!

Dolor et quisque habitant placerat bibendum risus; maecenas sem id
massa. Morbi vehicula sollicitudin rhoncus tincidunt, neque aliquet
tincidunt suspendisse. Pellentesque sociosqu mus scelerisque; bibendum
dis litora! Aliquam purus proin semper magna turpis; semper netus
vestibulum cras dapibus. Commodo\index{Commodo} condimentum inceptos
senectus.

Ipsum placerat pellentesque praesent lacus, convallis ultricies inceptos
fringilla lacinia tortor in integer! Aliquet fames penatibus justo
gravida, venenatis per ligula ullamcorper vestibulum fringilla sem. Non
vehicula curae pellentesque, dignissim, primis pellentesque tristique
accumsan dignissim tempus ridiculus. Turpis in porta eu vulputate.
Scelerisque leo mus: leo, penatibus tempor sodales inceptos maecenas.
Aptent est sodales in, sodales inceptos enim? Consequat sed erat odio
orci praesent morbi conubia: faucibus habitant parturient dapibus aenean
auctor lacus a!

Elit at massa tristique, enim bibendum suscipit et per libero. Duis
feugiat nullam morbi diam dignissim a commodo iaculis dictum rutrum! Vel
etiam libero: mauris lectus sollicitudin ultrices eu convallis ligula
nullam nisl blandit nullam.

\chapter{Introdução aos Indicadores Sociais e a Agenda
2030}\label{introduuxe7uxe3o-aos-indicadores-sociais-e-a-agenda-2030}

A análise de dados em ciências sociais e a formulação de políticas
públicas exigem métricas precisas e reprodutíveis. Este manual introduz
o desenvolvimento e a análise de indicadores sociais como recursos
fundamentais para o planejamento governamental e a atuação de
organizações não governamentais.

O acompanhamento rigoroso das condições de vida da população ganha
especial relevância no contexto da \textbf{Agenda 2030} e dos Objetivos
de Desenvolvimento Sustentável (ODS). O monitoramento local -- ao nível
municipal -- das metas da Agenda 2030 requer ferramentas analíticas
capazes de processar grandes volumes de dados de maneira transparente e
auditávegil.

\section{Validade, Fidedignidade e a Natureza dos
Indicadores}\label{validade-fidedignidade-e-a-natureza-dos-indicadores}

Um indicador social não é um dado bruto, mas uma construção metodológica
projetada para captar fenômenos sociais complexos, como pobreza,
educação e desigualdade. Para que um indicador seja útil à formulação de
políticas públicas, ele deve atender a dois critérios metodológicos
centrais:

\begin{enumerate}
\def\labelenumi{\arabic{enumi}.}
\tightlist
\item
  \textbf{Validade}: A capacidade do indicador de medir exatamente o
  fenômeno sociológico ou econômico a que se propõe.
\item
  \textbf{Fidedignidade (Confiabilidade)}: A consistência da métrica. O
  indicador deve apresentar estabilidade quando aplicado repetidas vezes
  sob as mesmas condições.
\end{enumerate}

Um dos exemplos mais consolidados dessa estruturação é o \textbf{Índice
de Desenvolvimento Humano (IDH)}. No Brasil, a adaptação metodológica
deste índice para a escala municipal (IDHM) é fruto da parceria entre o
Programa das Nações Unidas para o Desenvolvimento (PNUD), a Fundação
João Pinheiro (FJP) e o Instituto de Pesquisa Econômica Aplicada (Ipea),
materializada no Atlas Brasil. O Atlas disponibiliza mais de 200
indicadores sociais baseados nos dados dos Censos Demográficos do IBGE
(1991, 2000 e 2010).

Em termos formais, índices compostos frequentemente recorrem a
normalizações e médias ponderadas ou geométricas. A formulação
matemática clássica para agregação de subíndices de desenvolvimento pode
ser expressa em LaTeX da seguinte forma:

\[ 
IDH = \sqrt{I_{renda} \times I_{educacao} \times I_{longevidade}} 
\]

\chapter{O Paradigma Computacional: Por que usar o
R?}\label{o-paradigma-computacional-por-que-usar-o-r}

A transição de pacotes estatísticos clássicos para linguagens de
programação abertas representa um salto metodológico na pesquisa social.

\section{Interfaces de Apontar e Clicar vs.~Linguagens de
Script}\label{interfaces-de-apontar-e-clicar-vs.-linguagens-de-script}

Softwares tradicionais como o SPSS, desenvolvidos entre o final da
década de 1960 e o início dos anos 1970, popularizaram a análise de
dados por meio de interfaces de ``apontar e clicar''
(\emph{point-and-click}). Embora versáteis e fundamentais para a
história da pesquisa quantitativa, essas ferramentas limitam o
pesquisador às opções preestabelecidas nos menus do programa. Essa
limitação equivale a ``desenhar seguindo um tracejado'' imposto pelo
software.

Em contrapartida, o \textbf{R} é uma linguagem de programação
fundamentada na escrita de \emph{scripts} (roteiros de códigos
documentados). Desenvolvido inicialmente como um derivado de código
aberto da linguagem S (S-Plus), o R foi concebido para ser interativo. O
uso de \emph{scripts} libera o potencial analítico e criativo do
pesquisador, permitindo o desenvolvimento de novas rotinas, modelagens
complexas e, sobretudo, garantindo a \textbf{reprodutibilidade} da
pesquisa.

\begin{tcolorbox}[enhanced jigsaw, arc=.35mm, bottomrule=.15mm, bottomtitle=1mm, breakable, colback=white, colbacktitle=quarto-callout-important-color!10!white, colframe=quarto-callout-important-color-frame, coltitle=black, left=2mm, leftrule=.75mm, opacityback=0, opacitybacktitle=0.6, rightrule=.15mm, title=\textcolor{quarto-callout-important-color}{\faExclamation}\hspace{0.5em}{Reprodutibilidade como Princípio Ético}, titlerule=0mm, toprule=.15mm, toptitle=1mm]

Na análise de políticas públicas com dados do IBGE ou do Atlas Brasil, o
uso de \emph{scripts} em R garante que qualquer cidadão ou auditor possa
executar o mesmo código e chegar exatamente aos mesmos resultados,
configurando um pilar de transparência governamental.

\end{tcolorbox}

\section{A Comunidade e o Ecossistema Open
Source}\label{a-comunidade-e-o-ecossistema-open-source}

O R não é apenas um software, mas um ecossistema sustentado por uma
comunidade global altamente colaborativa. A arquitetura da linguagem
baseia-se em bibliotecas, geridas pela rede CRAN (\emph{Comprehensive R
Archive Network}). Atualmente, existem mais de 18.000 pacotes que
expandem as capacidades do programa para mineração de dados,
visualização gráfica, e aprendizado de máquina, rivalizando
positivamente com linguagens genéricas como o Python.

\chapter{Ambiente de Desenvolvimento: R e
RStudio}\label{ambiente-de-desenvolvimento-r-e-rstudio}

Para operar a linguagem, adotamos o \textbf{RStudio} (produto da empresa
\emph{Posit}), um Ambiente de Desenvolvimento Integrado (IDE) projetado
especificamente para maximizar a produtividade no R.

\section{Estrutura do RStudio}\label{estrutura-do-rstudio}

A interface do RStudio é classicamente dividida em quatro painéis: 1.
\textbf{Source / Editor de Script}: Área destinada à escrita, depuração
e documentação do código-fonte. 2. \textbf{Console}: O motor do R em si,
que aguarda os comandos interativos e exibe saídas numéricas e de texto.
3. \textbf{Environment}: O ambiente de trabalho, que lista as bases de
dados carregadas em memória, os objetos criados e o histórico de
execuções. 4. \textbf{Files / Plots / Packages / Help}: Painel
utilitário que gerencia o diretório de trabalho, exibe gráficos gerados,
lista pacotes e provê documentação.

\section{Processo de Instalação}\label{processo-de-instalauxe7uxe3o}

A configuração do ambiente exige a instalação sequencial de duas
plataformas:

\begin{enumerate}
\def\labelenumi{\arabic{enumi}.}
\tightlist
\item
  \textbf{R (Linguagem)}: Baixado através do repositório oficial CRAN
  (\texttt{https://cran.r-project.org}). Deve ser instalado antes da
  interface gráfica.
\item
  \textbf{RStudio (IDE)}: Disponível no site da Posit
  (\texttt{https://posit.co/download/rstudio-desktop/}). A versão
  Desktop \emph{Open Source} requer um sistema operacional de 64-bits
  (Windows, Mac ou Linux).
\end{enumerate}

\begin{tcolorbox}[enhanced jigsaw, arc=.35mm, bottomrule=.15mm, bottomtitle=1mm, breakable, colback=white, colbacktitle=quarto-callout-note-color!10!white, colframe=quarto-callout-note-color-frame, coltitle=black, left=2mm, leftrule=.75mm, opacityback=0, opacitybacktitle=0.6, rightrule=.15mm, title=\textcolor{quarto-callout-note-color}{\faInfo}\hspace{0.5em}{Compilação e RTools (Usuários de Windows)}, titlerule=0mm, toprule=.15mm, toptitle=1mm]

Alguns pacotes avançados (como certas dependências metodológicas) podem
requerer compilação do código-fonte, escrito em C/C++ ou Fortran. Para
usuários do sistema Windows, é necessário instalar adicionalmente o
\textbf{RTools} (disponível no CRAN). Este utilitário consome cerca de 1
GB de armazenamento. No entanto, rotinas básicas permitem ignorar o
aviso de compilação utilizando os arquivos binários já disponíveis.

\end{tcolorbox}

\chapter{Operações Básicas e Fluxos de Trabalho
Modernos}\label{operauxe7uxf5es-buxe1sicas-e-fluxos-de-trabalho-modernos}

Abaixo, apresentamos uma demonstração da sintaxe do \emph{Tidyverse}
moderno, utilizando o operador \emph{pipe} nativo do R
(\texttt{\textbar{}\textgreater{}}), para estruturar um processamento
hipotético dos dados do ``Atlas Brasil'' aplicados ao controle social.

\begin{Shaded}
\begin{Highlighting}[numbers=left,,]
\CommentTok{\# Carregamento de bibliotecas para manipulação de dados}
\FunctionTok{library}\NormalTok{(dplyr)}
\FunctionTok{library}\NormalTok{(readr)}

\CommentTok{\# Exemplo conceitual: Importação e filtragem de Indicadores Municipais}
\CommentTok{\# O foco recai na avaliação da agenda 2030 via IDHM (Atlas Brasil)}
\NormalTok{dados\_atlas\_brasil }\OtherTok{\textless{}{-}} \FunctionTok{read\_csv}\NormalTok{(}\StringTok{"caminho/para/dados\_atlas.csv"}\NormalTok{)}

\CommentTok{\# Processamento moderno utilizando o pipe nativo do R ( |\textgreater{} )}
\NormalTok{indicadores\_prioritarios }\OtherTok{\textless{}{-}}\NormalTok{ dados\_atlas\_brasil }\SpecialCharTok{|\textgreater{}}
  \CommentTok{\# Filtra municípios com dados consolidados para o ano do censo de 2010}
  \FunctionTok{filter}\NormalTok{(ano }\SpecialCharTok{==} \DecValTok{2010}\NormalTok{) }\SpecialCharTok{|\textgreater{}}
  \CommentTok{\# Seleciona variáveis de interesse socioeconômico}
  \FunctionTok{select}\NormalTok{(cod\_ibge, nome\_municipio, idhm, idhm\_renda, idhm\_educacao) }\SpecialCharTok{|\textgreater{}}
  \CommentTok{\# Ordena de forma decrescente pelo índice principal}
  \FunctionTok{arrange}\NormalTok{(}\FunctionTok{desc}\NormalTok{(idhm))}

\CommentTok{\# Visualização no console (limitação a 5 registros)}
\FunctionTok{head}\NormalTok{(indicadores\_prioritarios, }\DecValTok{5}\NormalTok{)}
\end{Highlighting}
\end{Shaded}

O uso metodológico de \emph{scripts} devidamente anotados não apenas
agiliza o trabalho do analista, mas garante que prefeituras e
pesquisadores possam reavaliar políticas sociais de forma consistente ao
longo do tempo.

\part{Referências}

\chapter*{References}\label{references}
\addcontentsline{toc}{chapter}{References}

\markboth{References}{References}

\protect\phantomsection\label{refs}

\chapter*{Sobre os autores}\label{sobre-os-autores}
\addcontentsline{toc}{chapter}{Sobre os autores}

\markboth{Sobre os autores}{Sobre os autores}

\begin{description}
\tightlist
\item[Ronaldo Baltar]
é sociólogo formado em 1985 pela Universidade de Brasília. No mesmo
período, concluiu curso técnico de programação de computadores e mantém
um aperfeiçoamento contínuo, sobretudo em linguagens e métodos de
análise de dados. Concluiu o mestrado em 1990 em Sociologia pela
Universidade Estadual de Campinas e em 1996, defendeu o Doutorado em
Sociologia, na Universidade de São Paulo. Desde então vem desenvolvendo
pesquisas, publicando artigos e ministrando cursos sobre análise
computacional de dados para a pesquisa aplicada em ciências sociais.
Atuou em diversos projetos de análise e avaliação de políticas públicas
e programas sociais. Atualmente é professor do Departamento de Ciências
Sociais da Universidade Estadual de Londrina, onde coordena o projeto
InfoSoc para formação complementar de graduação em informática aplicada
à pesquisa social, e participa como pesquisador do projeto Populações e
Dinâmica Social (\url{https://populacoes.info}).
\item[Cláudia Siqueira Baltar]
é demógrafa e professora do Departamento de Ciências Sociais da
Universidade Estadual de Londrina. Graduou-se em ciências Sociais em
2000 na Universidade de São Paulo. Concluiu o Mestrado em Ciência
Política na Universidade de Campinas em 2004 e defendeu tese de
Doutorado em 2008 no Programa de Pós-Graduação em Demografia da
Universidade de Campinas. Foi pesquisadora bolsista do Núcleo de Estudos
de População -- NEPO, na Universidade de Campinas. Atualmente é
coordenadora do projeto Populações e Dinâmica Social na Universidade
Estadual de Londrina (\url{https://populacoes.info}).
\end{description}


\backmatter


\end{document}
